% GENERATED FILE. Edit canonical lessons, then run npm run generate:books.
% canonical-source-hash: fnv1a64:bcae9bfc7d4e3463
% canonical-lessons: ES-C448-mudarse, ES-C448-mudanza, ES-C448-duro, ES-C448-estropear, ES-C448-roto, ES-R448-repaso-mudanza, ES-C448-sintesis-mudanza

\chapter{La mudanza}
\label{ch:es-mudanza}
\hlchaptermodality{448}

\begin{chapteropening}
\textbf{By the end of this chapter:} \emph{I can hand a flat over and take one on: say when I am moving, call the day hard, and tell somebody which things are out of order and which are actually in pieces.}

Read the message or notice this chapter's words exist for, then say the same thing back.
\end{chapteropening}

\section[mudarse]{mudarse — changing where you live}

\label{lesson:ES-C448-mudarse}



\begin{quote}
Say \textbf{la casa} and \textbf{el piso}. You are about to leave one of them.
\end{quote}

\subsection*{You'll want to know: mudarse}
\textbf{mudarse} — to move house.

\emph{Me mudo el sábado} — I move on Saturday. \emph{Nos mudamos a un piso más grande} — we are moving to a bigger flat. \emph{¿Cuándo te mudas?} — when do you move?

Spanish does not reach for a verb of carrying here. It reaches for a verb of \textbf{changing}, and that is the whole idea: you have not transported yourself, you have changed which address is yours.

\begin{grammarlens}[title={Grammar Lens: what the -se is doing}]
There is a plain verb \textbf{mudar} underneath, and it means to change something else:

\noindent
\begin{tabularx}{\linewidth}{@{}>{\raggedright\arraybackslash}X>{\raggedright\arraybackslash}X@{}}
\toprule
\textbf{without -se} & \textbf{with -se} \\
\midrule
\emph{la serpiente muda la piel} — sheds its skin & \emph{me mudo a Sevilla} — I move \\
\emph{mudar de opinión} — change your mind & \emph{se muda de ropa} — changes clothes \\
\bottomrule
\end{tabularx}

On the left the verb reaches out to something: a skin, an opinion. On the right it turns back and lands on the person. Once it lands on you, the thing being changed is understood to be \textbf{where you live} — or, in a bedroom, what you are wearing. Context decides between those two, and it decides easily, because nobody packs a lorry to change their shirt.

Note that \emph{mudarse} takes \textbf{a} for the destination — \emph{me mudo a Madrid} — the same \textbf{a} you already use for going anywhere.
\end{grammarlens}

\begin{cousinweb}
Latin \textbf{mutare}, to change. It is a wide family and you can hear it in English: a \textbf{mutation} is a change, to \textbf{mute} a sound is to alter it, and a bird in \textbf{moult} is changing its feathers — the same picture as the snake.

Spanish keeps the plain sense in \emph{mudar de opinión} and the specialised one in \emph{mudarse}. A word narrowing from \textquotedblleft{}change anything\textquotedblright{} to \textquotedblleft{}change address\textquotedblright{} is ordinary: the commonest thing people announce they are changing is where they sleep.
\end{cousinweb}

\subsection*{Your turn}
Say these aloud:

\begin{itemize}
\raggedright
  \item I am moving on Saturday
  \item we are moving to a bigger house
  \item what \emph{mudar} means without the \emph{-se}
\end{itemize}

\begin{tcolorbox}[breakable,colback=teal!4,colframe=teal!35!black,title={Before you move on}]
To move house? (\textbf{Mudarse}.) Which preposition for the destination? (\textbf{A} — \emph{me mudo a Bilbao}.) What Latin verb is inside it? (\textbf{Mutare}, to change.)
\end{tcolorbox}
\section[la mudanza]{la mudanza — the day itself}

\label{lesson:ES-C448-mudanza}



\begin{quote}
Say \textbf{mudarse} and \textbf{la caja}. Forty of the second, one of the first.
\end{quote}

\subsection*{You'll want to know: la mudanza}
\textbf{la mudanza} — the move: not the decision, the day.

\emph{Hacer una mudanza} — to do a move. \emph{La empresa de mudanzas} — the removal firm. \emph{El camión de mudanzas} — the removal van. \emph{La mudanza es el día quince} — the move is on the fifteenth.

Spanish says \textbf{hacer} a move where English says \emph{do} or \emph{make} one, and it never says \emph{tener}. A mudanza is something you carry out, not something you possess.

\begin{grammarlens}[title={Grammar Lens: -anza names the doing}]
Take a verb, cut the ending, add \textbf{-anza}, and you have a noun for the act or the state:

\noindent
\begin{tabularx}{\linewidth}{@{}>{\raggedright\arraybackslash}X>{\raggedright\arraybackslash}X@{}}
\toprule
\textbf{verb} & \textbf{noun} \\
\midrule
mudar & la mud\textbf{anza} — the move \\
esperar & la esper\textbf{anza} — hope \\
enseñar & la enseñ\textbf{anza} — teaching \\
tardar & la tard\textbf{anza} — the delay \\
\bottomrule
\end{tabularx}

All four are feminine, and so is every other \textbf{-anza} noun you will meet: \emph{la confianza}, \emph{la semejanza}, \emph{la balanza}. Spot the ending and the article comes free.

\emph{La esperanza} is worth a second look. From \emph{esperar}, to wait — hope, in Spanish, is named after the waiting rather than after the wanting. The ending took an ordinary verb and produced an abstract noun, and that is exactly what it does to \emph{mudar}, except that a mudanza is not abstract at all: it is boxes.
\end{grammarlens}

\begin{cousinweb}
Latin \textbf{mutare}, to change, plus the ending. Old Spanish used \emph{mudanza} for any change at all, including a change of heart, and you still meet that sense in older writing and in song lyrics.

The narrowing went the same way as the verb's. When people say \emph{la mudanza} now, boxes are in the room.
\end{cousinweb}

\subsection*{Your turn}
Say these aloud:

\begin{itemize}
\raggedright
  \item the move is on Saturday
  \item the removal firm is bringing the boxes
  \item which verb goes with \emph{mudanza}
\end{itemize}

\begin{tcolorbox}[breakable,colback=teal!4,colframe=teal!35!black,title={Before you move on}]
The move? (\textbf{La mudanza}.) The verb? (\textbf{Mudarse}.) And the gender that comes with \emph{-anza}? (\textbf{Feminine}.)
\end{tcolorbox}
\section[duro]{duro — hard, and hard going}

\label{lesson:ES-C448-duro}



\begin{quote}
Say \textbf{durante} and \textbf{fuerte}. One of them is about how long, the other about how much force.
\end{quote}

\subsection*{You'll want to know: duro}
\textbf{duro} — hard.

\emph{El pan está duro} — the bread has gone hard. \emph{Un colchón duro} — a firm mattress. And then, away from things you can press: \emph{un día muy duro} — a hard day. \emph{El trabajo es duro} — the work is heavy going. \emph{No seas tan duro con él} — don't be so hard on him.

A mudanza is the standard example of \emph{un día duro}, and no Spanish speaker would call it \emph{un día fuerte}.

\begin{grammarlens}[title={Grammar Lens: duro is not fuerte}]
The two get confused because English \textquotedblleft{}strong\textquotedblright{} and \textquotedblleft{}hard\textquotedblright{} overlap. Spanish keeps them apart by asking a different question about each:

\noindent
\begin{tabularx}{\linewidth}{@{}>{\raggedright\arraybackslash}X>{\raggedright\arraybackslash}X@{}}
\toprule
\textbf{word} & \textbf{the question it answers} \\
\midrule
\textbf{fuerte} & how much force does it have? \\
\textbf{duro} & how much does it resist? \\
\bottomrule
\end{tabularx}

\emph{Una persona fuerte} can lift the boxes. \emph{Una persona dura} will not soften when you ask for a favour. A rock is \emph{duro}; it has no force at all. Wind is \emph{fuerte}; it is not hard in the least.

So when a day is \emph{duro}, the picture is not that the day was powerful. It is that the day would not give.
\end{grammarlens}

\begin{cousinweb}
Latin \textbf{durus}, hard — and from it the verb \textbf{durare}, to last.

That derivation repays a moment. What lasts is what resists wearing away, so \textquotedblleft{}hard\textquotedblright{} and \textquotedblleft{}long-lasting\textquotedblright{} are one idea seen twice. Spanish kept both halves: \emph{duro} is the resisting, and \textbf{durar} — with \textbf{durante} inside it — is the lasting. \emph{La mudanza duró todo el día}: the move held out all day.

English took both halves too, from the same place: \textbf{durable} and \textbf{duration} on one side, \textbf{endure} on the other.
\end{cousinweb}

\subsection*{Your turn}
Say these aloud:

\begin{itemize}
\raggedright
  \item the bread is hard
  \item the move was a very hard day
  \item why \emph{durante} and \emph{duro} are relatives
\end{itemize}

\begin{tcolorbox}[breakable,colback=teal!4,colframe=teal!35!black,title={Before you move on}]
Hard? (\textbf{Duro}.) Which word would you pick for a strong wind? (\textbf{Fuerte} — force, not resistance.) What links \emph{duro} to \emph{durar}? (\textbf{What resists is what lasts}.)
\end{tcolorbox}
\section[estropear / estropeado]{estropear / estropeado — it no longer works}

\label{lesson:ES-C448-estropear}



\begin{quote}
Say \textbf{torpe} and \textbf{romper}. Being the first often leads to the second.
\end{quote}

\subsection*{You'll want to know: estropear}
\textbf{estropear} — to ruin, to wreck, to put out of action.

\emph{He estropeado el móvil} — I've wrecked my phone. \emph{La lluvia estropeó la fiesta} — the rain ruined the party. With \textbf{-se} it happens by itself: \emph{se ha estropeado la lavadora} — the washing machine has broken down.

And the form you meet most is the adjective:

\textbf{estropeado} — out of order, damaged, spoiled.

\begin{quote}
Ascensor \textbf{estropeado}. Usen las escaleras.
\end{quote}

That is a real sign, on a real lift, and it is the reason this word earns its place before the polite ones.

\begin{grammarlens}[title={Grammar Lens: estropeado is not roto}]
Both translate as \textquotedblleft{}broken\textquotedblright{} and they are not interchangeable:

\noindent
\begin{tabularx}{\linewidth}{@{}>{\raggedright\arraybackslash}X>{\raggedright\arraybackslash}X@{}}
\toprule
\textbf{what you mean} & \textbf{the word} \\
\midrule
it is in pieces & \textbf{roto} \\
it will not work & \textbf{estropeado} \\
\bottomrule
\end{tabularx}

A washing machine full of water that does nothing is \textbf{estropeada} and not \emph{rota} — nothing came apart. A plate in two halves is \textbf{rota} and calling it \emph{estropeada} would sound like an understatement.

Food takes \emph{estropeado} too: \emph{la fruta se ha estropeado} — the fruit has gone off. Again the test holds. The fruit is not in pieces; it no longer does its job.
\end{grammarlens}

\begin{cousinweb}
Spanish took it from Italian \textbf{stroppiare}, to cripple or maim, which goes back to Latin \textbf{turpis} — shameful, ugly.

So \textbf{estropear} and \textbf{torpe} are cousins, and the pair is worth holding side by side. The same Latin word for \emph{shameful} came out of Italian as \textquotedblleft{}wreck a thing\textquotedblright{} and out of Spanish as \textquotedblleft{}clumsy\textquotedblright{} — an apology you make over a spilled glass. One branch kept the weight and one branch lost nearly all of it.

The \textbf{e-} at the front is Spanish's own doing. Spanish will not open a word with \emph{st-}, so it fits a vowel in front: \emph{stroppiare} became \emph{estropear}, the way \emph{stare} became \emph{estar} and \emph{studiare} became \emph{estudiar}.
\end{cousinweb}

\subsection*{Your turn}
Say these aloud:

\begin{itemize}
\raggedright
  \item the lift is out of order
  \item the washing machine has broken down
  \item which of \emph{roto} and \emph{estropeado} fits a plate in two halves
\end{itemize}

\begin{tcolorbox}[breakable,colback=teal!4,colframe=teal!35!black,title={Before you move on}]
Out of order? (\textbf{Estropeado}.) The verb? (\textbf{Estropear}.) Which word already in your head shares its Latin root? (\textbf{Torpe} — both from \emph{turpis}.)
\end{tcolorbox}
\section[roto]{roto — in pieces}

\label{lesson:ES-C448-roto}



\begin{quote}
Say \textbf{romper} and \textbf{la caja}. Now say what happened to the box.
\end{quote}

\subsection*{You'll want to know: roto}
\textbf{roto} — broken, torn, in pieces.

\emph{Un vaso roto} — a broken glass. \emph{Los pantalones están rotos} — the trousers are torn. \emph{He roto el plato} — I've broken the plate.

It agrees like any adjective: \emph{roto, rota, rotos, rotas}.

\begin{grammarlens}[title={Grammar Lens: the participle that lost its middle}]
A regular \emph{-er} verb would give \emph{romp-ido}. Spanish says \textbf{roto}, and it is one of a small group that behave this way:

\noindent
\begin{tabularx}{\linewidth}{@{}>{\raggedright\arraybackslash}X>{\raggedright\arraybackslash}X>{\raggedright\arraybackslash}X@{}}
\toprule
\textbf{verb} & \textbf{the participle you expect} & \textbf{the one Spanish uses} \\
\midrule
romper & \emph{rompido} & \textbf{roto} \\
escribir & \emph{escribido} & \textbf{escrito} \\
abrir & \emph{abrido} & \textbf{abierto} \\
\bottomrule
\end{tabularx}

These short forms are the older ones. They came straight down from Latin participles that were already short, and because the verbs are used so constantly, the regular pattern never got a grip on them. The same thing happens in English, where \emph{wrote} survives beside \emph{worked}.

One form covers both jobs: \emph{he roto el plato} after \emph{haber}, and \emph{el plato está roto} as a plain description.
\end{grammarlens}

\begin{cousinweb}
Latin \textbf{rumpere}, to burst, with participle \textbf{ruptus} — and \emph{ruptus} is what became \emph{roto}.

English borrowed the same participle whole: a \textbf{rupture}, an \textbf{eruption}, something \textbf{abrupt} that breaks off. Spanish let it wear down into a two-syllable everyday word, and kept \emph{ruptura} alongside for the borrowed, formal sense.

Keep the pair from the lesson before this one in view. \emph{Roto} is \emph{ruptus}, bursting. \emph{Estropeado} is \emph{turpis}, shameful. Spanish named the two kinds of failure after two different pictures, and the pictures still sort them.
\end{cousinweb}

\subsection*{Your turn}
Say these aloud:

\begin{itemize}
\raggedright
  \item a broken glass, then broken trousers
  \item I have broken the plate
  \item whether a dead fridge is \emph{roto} or \emph{estropeado}
\end{itemize}

\begin{tcolorbox}[breakable,colback=teal!4,colframe=teal!35!black,title={Before you move on}]
Broken, in pieces? (\textbf{Roto}.) What would the regular participle have been? (\emph{Rompido}.) Name two more verbs with a short participle. (\textbf{Escrito}, \textbf{abierto}.)
\end{tcolorbox}
\section[(mudarse, la mudanza, duro …]{(mudarse, la mudanza, duro, estropeado, roto) — from cold}

\label{lesson:ES-R448-repaso-mudanza}



\begin{quote}
Close the lessons before this one. Say all five: \emph{to move house}, \emph{the move}, \emph{hard}, \emph{out of order}, \emph{in pieces}.
\end{quote}

\begin{grammarlens}[title={Grammar Lens: two words for broken, one question}]
This is the one distinction in the chapter that will cost you marks if it goes unsorted. Ask: \textbf{did anything come apart?}

\noindent
\begin{tabularx}{\linewidth}{@{}>{\raggedright\arraybackslash}X>{\raggedright\arraybackslash}X@{}}
\toprule
\textbf{the situation} & \textbf{the word} \\
\midrule
lift stuck between floors & \textbf{estropeado} \\
window with a hole in it & \textbf{roto} \\
milk that has gone off & \textbf{estropeado} \\
shoelace snapped & \textbf{roto} \\
\bottomrule
\end{tabularx}

Notice the milk. Nothing could come apart in a carton, so \emph{roto} is not available — and that is the test doing its work on a case where English gives you no help at all, since English would say neither \textquotedblleft{}broken milk\textquotedblright{} nor \textquotedblleft{}out-of-order milk\textquotedblright{}.
\end{grammarlens}

\begin{grammarlens}[title={Grammar Lens: where the five came from}]
Four of the five were built out of things you already had.

\noindent
\begin{tabularx}{\linewidth}{@{}>{\raggedright\arraybackslash}X>{\raggedright\arraybackslash}X@{}}
\toprule
\textbf{new word} & \textbf{what it was built on} \\
\midrule
\textbf{mudarse} & \emph{mudar}, to change \\
\textbf{la mudanza} & \emph{mudar} plus the \emph{-anza} of \emph{esperanza} \\
\textbf{duro} & the \emph{dur-} inside \emph{durante} \\
\textbf{estropeado} & \emph{turpis}, the Latin behind \emph{torpe} \\
\bottomrule
\end{tabularx}

Only \textbf{roto} arrived on its own, and even it came with \emph{romper} attached. Five words, one genuinely new root among them. That is what happens once a vocabulary is large enough: new words start arriving as relatives.
\end{grammarlens}

\subsection*{Your turn}
Say these aloud:

\begin{itemize}
\raggedright
  \item we are moving on the fifteenth
  \item the move was a hard day
  \item which preposition follows \emph{mudarse}
\end{itemize}

\begin{tcolorbox}[breakable,colback=teal!4,colframe=teal!35!black,title={Before you move on}]
All five again, in Spanish, without looking. Then one sentence using two of them at once.
\end{tcolorbox}
\section[Practice]{(el día de la mudanza) — read it, then do it yourself}

\label{lesson:ES-C448-sintesis-mudanza}



\begin{quote}
Say \textbf{la mudanza} and \textbf{estropeado}. Both are in the note below, and one of them is doing the work of a warning.
\end{quote}

\subsection*{Your turn}
A note left on the kitchen counter for whoever moves in next. Read it, then answer.

\begin{quote}
Hola. Nos \textbf{mudamos} hoy y te dejo unas cosas por escrito.
\end{quote}

\begin{quote}
El timbre está \textbf{estropeado} desde enero. El portero dice que lo arregla, pero lleva así medio año, así que llama por teléfono.
\end{quote}

\begin{quote}
En el armario hay un cristal \textbf{roto}. Lo he tapado con cartón, pero ten cuidado al abrirlo.
\end{quote}

\begin{quote}
La \textbf{mudanza} ha sido \textbf{dura} y quedan tres cajas en el trastero. Las recojo el domingo.
\end{quote}

Say these aloud:

\begin{itemize}
\raggedright
  \item which two things in the flat are not in working order
  \item which of the two came apart, and how you know
  \item what is still in the storage room, and when it goes
\end{itemize}

\begin{culture}
The note uses both words for broken within four lines, and each one is the only correct choice.

The doorbell is \textbf{estropeado}. Nothing about it is in pieces; it has stopped doing its job, and the fault is invisible. The pane of glass is \textbf{roto}, which is why it needed covering with cardboard — there is now a hole where glass used to be. Swap the two words and a Spanish reader would picture a smashed doorbell and a window that mysteriously fails to be a window.

Then look at \textbf{nos mudamos hoy}. Present tense, for something happening today, exactly as in English. And \textbf{dura}, agreeing with \emph{la mudanza}: the day resisted, all day long.
\end{culture}

\begin{tcolorbox}[breakable,colback=teal!4,colframe=teal!35!black,title={Before you move on}]
Now your turn, out loud and then in writing. Leave a note of four lines for the next tenant: say you are moving today, name one thing that is out of order, name one thing that is broken, and say the move has been hard.

Then check yourself: did each of \emph{roto} and \emph{estropeado} go to the right object, and did \emph{duro} agree with what it describes?
\end{tcolorbox}
